This is the first of two marked assignments for the course ``Proof Assistants''
(L81), covering the Isabelle part of the course.  This assignment comprises
a number of small formalisation and verification tasks related to the imperative
programming language IMP, as well as a short write-up explaining your work.

The due date for this assignment is Thursday, 5 March 2026, 4pm. You must
work on this assignment as an individual;  collaboration with others or the use
of AI tools (other than the Sledgehammer tool built into Isabelle) is not
allowed. Copying material found elsewhere counts as plagiarism. Please use the
Moodle page for the course to submit an archive containing your theories and
write-up by the deadline.

Each of the two assignments is worth 100 marks, distributed as follows:
\begin{itemize}
  \item 50 marks for completing basic formalisation and verification tasks
    assessing grasp of the material taught in the lecture.  For this
    assignment, please use the Isabelle theory file template provided on the
    course website. When preparing your submission, please include your
    \verb|Assignment1.thy| file as well as any other theory files that you have
    edited.  You are allowed to use Sledgehammer to find proofs, although you
    might want to streamline any complex proofs it finds into something more
    readable.  The main assessment criteria for these tasks are correctness and
    completeness of the specifications and proofs.  You might want to test your
    specifications to check that they work as intended.
  \item 20 marks for completing designated more challenging tasks, which might
    require more creativity when designing specifications and proof strategies,
    and might require some Isabelle techniques that go beyond what was taught
    in the lecture. Completing these advanced tasks can help you achieve a
    distinction grade, but you might want to focus first on completing the main
    tasks before attempting the advanced tasks.
  \item 30 marks for a clear write-up, where 10 of these marks will be reserved
    for write-ups of exceptional quality, e.g. demonstrating particular
    insight. In general, your write-up should explain the design decisions you
    made during the formalisation and the strategies you used for the proofs.
    It might also discuss proof attempts that failed in interesting ways and
    lessons you learned from them, if that happens. The maximum length of the
    write-up is 2,500 words, although it could be much shorter, especially if
    you add comments to your theory files. It is possible to use Isabelle to
    generate a document,\footnote{The template files for the assignment are set
    up to generate the handout document with this introduction and the tasks.
    For more information on document preparation from Isabelle, see Chapter 4
    of ``The Isabelle/Isar Reference Manual'' and Chapter 3 of ``The Isabelle
    System Manual''.} but you can use any tool you prefer to prepare your
    write-up and add a PDF to your submission.
\end{itemize}
